\documentclass[11pt]{article}
\usepackage[a4paper,margin=22mm]{geometry}
\usepackage{fontspec}
\setmainfont{DejaVu Serif}
\setsansfont{DejaVu Sans}
\usepackage{microtype}
\usepackage{xcolor}
\usepackage{hyperref}
\usepackage{enumitem}
\usepackage{parskip}
\definecolor{Accent}{HTML}{971D32}
\definecolor{Muted}{HTML}{666666}
\hypersetup{colorlinks=true,urlcolor=Accent,linkcolor=Accent}
\setlist{leftmargin=1.6em,itemsep=0.3em,topsep=0.35em}
\setcounter{tocdepth}{2}
\renewcommand{\contentsname}{Contents}
\newcommand{\sectionrule}{\par\vspace{-0.5em}\noindent\textcolor{Accent}{\rule{\linewidth}{0.6pt}}\par}
\begin{document}

\begin{center}
{\sffamily\bfseries\color{Accent} TOEFL WORD OF THE DAY\par}
\vspace{8mm}
{\fontsize{42}{48}\selectfont\bfseries determine\par}
\vspace{2mm}
{\Large\sffamily\color{Accent} /dɪˈtɝːmɪn/\par}
\vspace{2mm}
{\small\sffamily Local audio phrase: determine\par}
{\small\sffamily Audio file: audios/determine.mp3\par}
\vspace{2mm}
{\large\sffamily\color{Muted} transitive verb\par}
\vspace{5mm}
{\sffamily discover or establish \textbullet\ control an outcome \textbullet\ make an official decision\par}
\vspace{6mm}
{\large Determine can mean finding an answer, shaping an outcome, or making a formal decision.\par}
\vspace{6mm}
\begin{minipage}{0.84\textwidth}
\itshape “Researchers are trying to determine whether sleep quality affects long-term memory.”
\end{minipage}\par
\vspace{10mm}
\textbf{Date:} September 16th 2026\quad\textbullet\quad \textbf{Level:} B1--mid-B2 TOEFL
\vfill
Created and compiled by \textbf{Amir Kooshky}\par
\vspace{2mm}
\href{https://t.me/KooshkyTOEFL}{Telegram Channel}\quad\textbullet\quad
\href{https://www.instagram.com/kooshkytoefl}{Instagram}\quad\textbullet\quad
\href{https://t.me/KooshkyTOEFL\_pv}{Personal Telegram}
\end{center}
\newpage
\tableofcontents
\newpage

\section{Meanings and usage}
\sectionrule
\subsection{Meaning 1: Discover or establish}
\textbf{Part of speech:} transitive verb

\textbf{IPA:} /dɪˈtɝːmɪn/

\textbf{Pronunciation phrase:} determine

\textbf{Local audio file:} audios/determine.mp3

\textbf{Register:} neutral to formal; very common in academic and technical writing

\textbf{Definition:} to discover, calculate, or establish a fact, amount, cause, or answer through careful examination

\textbf{In simpler words:} find out or establish

\textbf{Typical pattern:} determine + noun, whether or if, question word, or that-clause

\subsubsection{Natural examples}
\begin{enumerate}
  \item Researchers are trying to determine whether sleep quality affects long-term memory.
  \item The laboratory test can determine the concentration of lead in the water.
  \item Investigators could not determine the exact cause of the equipment failure.
  \item Scientists examined the fossils to determine their approximate age.
  \item The survey was designed to determine how residents use public transportation.
  \item Further evidence is needed to determine which explanation is more accurate.
  \item The committee will determine whether the proposal meets the legal requirements.
  \item Analysts used satellite images to determine the extent of forest loss.
\end{enumerate}

\subsubsection{Nuance and implication}
\textbf{Central nuance:} Determine suggests reaching a sufficiently definite answer by using evidence, measurement, calculation, or authority.

\textbf{Usual implication:} The answer was previously uncertain and requires investigation or judgment.

\textbf{Compared with estimate:} To estimate is to produce an approximate value. To determine is to establish an answer considered sufficiently definite for the purpose.

\subsubsection{Grammar notes}
\textbf{How to use it:} This verb normally needs a fact or question after it.

\textbf{What can follow it:} Use a noun, as in determine the cause; whether or if, as in determine whether it works; or a question word, as in determine how it happened.

\textbf{Position in a sentence:} It is frequent in passive forms such as The age was determined by carbon dating.

\subsubsection{Useful patterns}
\textbf{Pattern:} determine + cause, effect, amount, extent, value, or identity

\textbf{Use:} Use this when an investigation establishes a fact.

\textbf{Example:} The study sought to determine the extent of the damage.

\textbf{Pattern:} determine whether or if + clause

\textbf{Use:} Use this when choosing between possible answers.

\textbf{Example:} Additional trials will determine whether the treatment is safe.

\textbf{Pattern:} determine how, why, when, or where + clause

\textbf{Use:} Use a question word to identify the information being investigated.

\textbf{Example:} Researchers hope to determine why the population declined.

\subsubsection{Common wrong usage}
\paragraph{Correction 1}
\textbf{Wrong:} The test determined about the cause of the infection.

\textbf{Correct:} The test determined the cause of the infection.

\textbf{Why:} Do not use about before a direct object such as the cause.

\paragraph{Correction 2}
\textbf{Wrong:} Scientists determined that how the material changed.

\textbf{Correct:} Scientists determined how the material changed.

\textbf{Why:} Use either that or a question word such as how, not both together.

\subsection{Meaning 2: Control or shape an outcome}
\textbf{Part of speech:} transitive verb

\textbf{IPA:} /dɪˈtɝːmɪn/

\textbf{Pronunciation phrase:} determine

\textbf{Local audio file:} audios/determine.mp3

\textbf{Register:} neutral; common in academic explanations of causes and outcomes

\textbf{Definition:} to control, cause, or strongly influence the nature or result of something

\textbf{In simpler words:} control or decide what happens

\textbf{Typical pattern:} factor + determines + outcome or whether/how clause

\subsubsection{Natural examples}
\begin{enumerate}
  \item Access to reliable transportation can determine whether workers accept jobs far from home.
  \item Soil quality largely determines which crops can grow in a region.
  \item A society's institutions help determine how resources are distributed.
  \item Early experiences do not completely determine a person's future behavior.
  \item The availability of water may determine the size of the settlement.
  \item Several interacting genes determine the color of the flowers.
  \item The final examination will determine half of each student's grade.
  \item Economic conditions often determine the pace of construction.
\end{enumerate}

\subsubsection{Nuance and implication}
\textbf{Central nuance:} In this sense, determine means having a decisive or powerful effect, not discovering information.

\textbf{Usual implication:} The subject is an important causal or controlling factor, although context may show that other factors also matter.

\textbf{Compared with influence:} Influence means affect to some degree. Determine often suggests a stronger effect that fixes or decides the outcome.

\subsubsection{Grammar notes}
\textbf{How to use it:} The controlling factor is the subject, and the result is placed after determine.

\textbf{What can follow it:} Use a noun or a clause beginning with whether, how, what, or which.

\textbf{Position in a sentence:} Adverbs such as largely, partly, ultimately, and directly can show how strong the effect is.

\subsubsection{Useful patterns}
\textbf{Pattern:} factor + determine(s) + outcome

\textbf{Use:} Use this to connect a controlling cause with its result.

\textbf{Example:} Temperature determines the rate of the chemical reaction.

\textbf{Pattern:} determine whether or how + clause

\textbf{Use:} Use this when a factor controls whether or how something happens.

\textbf{Example:} Household income may determine whether a student can remain in school.

\textbf{Pattern:} be determined by + factor

\textbf{Use:} Use the passive when the result is more important than the controlling factor.

\textbf{Example:} The final cost is determined by the amount of material used.

\subsubsection{Common wrong usage}
\paragraph{Correction 1}
\textbf{Wrong:} Weather determines on the success of the harvest.

\textbf{Correct:} Weather determines the success of the harvest.

\textbf{Why:} Determine takes a direct object; do not add on.

\paragraph{Correction 2}
\textbf{Wrong:} Income is determined from several social factors.

\textbf{Correct:} Income is determined by several social factors.

\textbf{Why:} Use by to introduce the controlling factor in a passive sentence.

\subsection{Meaning 3: Make an official decision}
\textbf{Part of speech:} transitive verb

\textbf{IPA:} /dɪˈtɝːmɪn/

\textbf{Pronunciation phrase:} determine

\textbf{Local audio file:} audios/determine.mp3

\textbf{Register:} formal when it means make an official decision

\textbf{Definition:} to make an official or authoritative decision about something after considering the facts

\textbf{In simpler words:} decide officially

\textbf{Typical pattern:} determine + noun, question-word clause, or that-clause

\subsubsection{Natural examples}
\begin{enumerate}
  \item After reviewing the evidence, the board determined that the investigation should continue.
  \item The court will determine the appropriate level of compensation.
  \item Officials determined that the building should remain closed until repairs were complete.
  \item The committee is responsible for determining who receives the scholarship.
  \item The agency determined the final boundaries of the protected area.
  \item The judge will determine the appropriate penalty after hearing both sides.
  \item The panel determined that a second vote was necessary.
  \item Local authorities determine the standards that new buildings must meet.
\end{enumerate}

\subsubsection{Nuance and implication}
\textbf{Central nuance:} Determine is more formal than decide here and often suggests that a court, committee, agency, or rule has authority to settle the matter.

\textbf{Usual implication:} The decision is reached after facts, standards, or evidence have been considered.

\textbf{Compared with decide:} Decide is the ordinary general word. Determine is more formal and is common when an authority settles a question.

\subsubsection{Grammar notes}
\textbf{How to use it:} Use determine followed by a noun, question-word clause, or that-clause.

\textbf{What can follow it:} Common forms include determine the outcome, determine who qualifies, and determine that action is needed.

\textbf{Position in a sentence:} Do not confuse this verb with the adjective determined in be determined to.

\subsubsection{Useful patterns}
\textbf{Pattern:} determine + question word + clause

\textbf{Use:} Use this for an official decision about a person, time, place, or method.

\textbf{Example:} The panel will determine who is eligible for funding.

\textbf{Pattern:} determine that + clause

\textbf{Use:} Use this to report a formal conclusion or decision.

\textbf{Example:} The court determined that the regulation was valid.

\textbf{Pattern:} determine + noun

\textbf{Use:} Use this when an authority decides an outcome, amount, boundary, or standard.

\textbf{Example:} A national panel determines the minimum safety standard.

\subsubsection{Common wrong usage}
\paragraph{Correction 1}
\textbf{Wrong:} The panel determined about the final amount.

\textbf{Correct:} The panel determined the final amount.

\textbf{Why:} Determine takes the matter being decided as a direct object.

\paragraph{Correction 2}
\textbf{Wrong:} The panel determined about who would receive the award.

\textbf{Correct:} The panel determined who would receive the award.

\textbf{Why:} A question-word clause follows determine directly.

\section{Word family}
\sectionrule
\subsection{determination}
\textbf{Part of speech:} countable or uncountable noun

\textbf{IPA:} /dɪˌtɝːməˈneɪʃən/

\textbf{Definition:} the process of establishing something, an official decision, or the strong quality of continuing toward a goal

\textbf{Relationship to the headword:} This noun covers both establishing an answer and being firmly resolved.

\textbf{Grammar pattern:} determination of + fact, or determination to + verb

\textbf{Usage note:} Its exact meaning depends strongly on context.

\subsubsection{Examples}
\begin{enumerate}
  \item The determination of the sample's age required further testing.
  \item Her determination helped her complete the demanding program.
\end{enumerate}

\subsection{determined}
\textbf{Part of speech:} adjective

\textbf{IPA:} /dɪˈtɝːmɪnd/

\textbf{Definition:} having made a firm decision and refusing to give up

\textbf{Relationship to the headword:} This adjective describes firm intention or effort.

\textbf{Grammar pattern:} be determined to + base verb, or a determined + effort

\textbf{Usage note:} Pronounce the final syllable as /mɪnd/, not like the verb mind.

\subsubsection{Examples}
\begin{enumerate}
  \item She was determined to improve access to rural schools.
  \item The team made a determined effort to recover the missing data.
\end{enumerate}

\section{Common collocations}
\sectionrule
\subsection{determine the cause}
\textbf{Applies to:} Discover or establish

\textbf{Meaning:} establish why something happened

\textbf{Example:} Investigators collected samples to determine the cause of the contamination.

\subsection{determine the extent}
\textbf{Applies to:} Discover or establish

\textbf{Meaning:} establish how large or serious something is

\textbf{Example:} More surveys are needed to determine the extent of the decline.

\subsection{largely determine}
\textbf{Applies to:} Control or shape an outcome

\textbf{Meaning:} have the main controlling influence on something

\textbf{Example:} Rainfall patterns largely determine which plants survive.

\subsection{be determined by}
\textbf{Applies to:} Control or shape an outcome

\textbf{Meaning:} be controlled or decided by a factor

\textbf{Example:} Eligibility is determined by household income.

\subsection{determine eligibility}
\textbf{Applies to:} Make an official decision

\textbf{Meaning:} decide officially who qualifies

\textbf{Example:} The committee uses published criteria to determine eligibility.

\section{Synonyms and antonyms}
\sectionrule
\subsection{Close synonyms}
\subsubsection{establish}
\textbf{Part of speech:} verb

\textbf{Applies to:} Discover or establish

\textbf{Shared meaning:} Both can mean reaching a firm factual conclusion.

\textbf{Difference:} Establish emphasizes showing that a fact is true. Determine can also involve calculating or identifying an answer.

\textbf{Example:} The evidence established that the document was authentic.

\subsubsection{decide}
\textbf{Part of speech:} verb

\textbf{Applies to:} Make an official decision

\textbf{Shared meaning:} Both can refer to choosing an outcome after consideration.

\textbf{Difference:} Decide is the usual everyday term. Determine is more formal when an authority reaches a decision.

\textbf{Example:} The committee decided to postpone the vote.

\section{Words learners may confuse}
\sectionrule
\subsection{determine}
\textbf{Part of speech:} verb versus adjective

\textbf{Why learners confuse them:} The forms look similar but perform different grammatical jobs.

\textbf{Key difference:} Determine is a verb. Determined is usually an adjective when it describes firm intention.

\textbf{determine:} The test will determine the cause.

\textbf{determine:} The researcher was determined to find the cause.

\section{Etymology}
\sectionrule
\textbf{Origin:} Determine comes from Latin determinare, meaning to limit or set boundaries.

\textbf{Development:} The idea of setting limits developed into fixing an answer, controlling a result, or settling a decision.

\textbf{Memory link:} To determine something is to bring uncertainty to a defined endpoint.

\section{Practice exercises}
\sectionrule
\subsection{Word family choice}
Choose the best member of the word family for each blank.

\begin{enumerate}
  \item The laboratory must \_\_\_\_\_ the chemical composition of the sample.
  \item Her \_\_\_\_\_ to finish the research never weakened.
  \item The team was \_\_\_\_\_ to identify the source of the error.
\end{enumerate}

\subsection{Correct or incorrect?}
Decide whether each sentence is correct. Correct the inaccurate ones.

\begin{enumerate}
  \item The survey aims to determine whether attitudes have changed.
  \item Climate partly determines which crops can grow in a region.
  \item She determined to completing the study.
  \item The cost is determined from the materials used.
\end{enumerate}

\subsection{Collocation practice}
Complete each sentence with a natural collocation.

\begin{enumerate}
  \item Tests were conducted to determine the \_\_\_\_\_ of the structural failure.
  \item The final price is determined \_\_\_\_\_ the amount of energy used.
\end{enumerate}

\subsection{Complete answer key}
\subsubsection{Word family choice}
\begin{enumerate}
  \item \textbf{determine.} The laboratory must determine the chemical composition of the sample. \emph{Use the base verb after must.}
  \item \textbf{determination.} Her determination to finish the research never weakened. \emph{Use the noun determination as the subject.}
  \item \textbf{determined.} The team was determined to identify the source of the error. \emph{Use be determined to for firm intention.}
\end{enumerate}

\subsubsection{Correct or incorrect?}
\begin{enumerate}
  \item \textbf{Correct} \emph{Determine whether correctly introduces the question being investigated.}
  \item \textbf{Correct} \emph{Here determine means strongly influence or control an outcome.}
  \item \textbf{Incorrect}. She was determined to complete the study. \emph{Use be determined to followed by the base verb for firm intention.}
  \item \textbf{Incorrect}. The cost is determined by the materials used. \emph{Use by for the controlling factor in a passive sentence.}
\end{enumerate}

\subsubsection{Collocation practice}
\begin{enumerate}
  \item \textbf{cause.} Tests were conducted to determine the cause of the structural failure. \emph{Determine the cause is a common academic collocation.}
  \item \textbf{by.} The final price is determined by the amount of energy used. \emph{Use determined by to name a controlling factor.}
\end{enumerate}

\vfill
\begin{center}
\small\color{Muted} Created and compiled by \textbf{Amir Kooshky}\\
\href{https://t.me/KooshkyTOEFL}{Telegram Channel}\quad\textbullet\quad
\href{https://www.instagram.com/kooshkytoefl}{Instagram}\quad\textbullet\quad
\href{https://t.me/KooshkyTOEFL\_pv}{Personal Telegram}
\end{center}
\end{document}
